\input qpd-bootstrap

\begin{problem}{20260731}
  \meta{name}{Swing Dynamics}
  \meta{setter}{Jason Zhang}
  \meta{score}{10}
  \meta{topic}{mechanics}
  \meta{difficulty}{5}
  \meta{tags}{mechanics, pendulum, circular motion, energy, centripetal force}

A physics study group builds a small model of a swing.  A point mass $M$ hangs from a fixed pivot on a light string of length $L$.  Let $\theta$ be the angle the swing makes with the downward vertical.

At $t = 0$ the swing is at $\theta = 0$, but the mass has a velocity of magnitude $v_0$ directed normal to the string.  At some later time $t'$, the tension in the string returns to zero.

\begin{assumptions}
  \item The string is massless and inextensible, of length $L$.
  \item The pivot is fixed; the only force besides tension is gravity $g$.
  \item The mass is a point particle; air resistance is negligible.
\end{assumptions}

\begin{questions}
  \item \textbf{Over the horizontal.}
        Derive the minimum $v_0$ such that at $t = t'$, $\theta > \pi/2$ (the swing has passed the horizontal).

  \item \textbf{Over the top.}
        Derive the minimum $v_0$ such that at $t = t'$, $\theta > \pi$ (the swing has passed the vertical).

  \item \textbf{The point of no return.}
        Find the minimum angle $\theta^{*}$ with the following property: if the swing reaches $\theta = \theta^{*}$ with the string still taut, then the tension is positive at every angle beyond $\theta^{*}$.  (In other words, $\theta^{*}$ is the point of no return — once passed with the string taut, the swing is guaranteed to remain taut for the rest of the revolution.)
\end{questions}

\end{problem}

\begin{solution}{20260731}

Let $\theta$ be measured from the downward vertical.  The radial (outward) equation of motion is
\[
  T - Mg\cos\theta = -\frac{Mv^{2}}{L}
  \qquad\Longrightarrow\qquad
  T = M\frac{v^{2}}{L} + Mg\cos\theta,
\]
where $T \ge 0$ as long as the string is taut.  Conservation of energy gives
\[
  \tfrac12 M v^{2} + MgL(1-\cos\theta) = \tfrac12 M v_{0}^{2}
  \qquad\Longrightarrow\qquad
  v^{2} = v_{0}^{2} - 2gL(1-\cos\theta).
\]

\subsection*{Q1. Over the horizontal}

The string goes slack when $T = 0$:
\[
  \frac{v^{2}}{L} + g\cos\theta = 0
  \qquad\Longrightarrow\qquad
  v^{2} = -gL\cos\theta .
\]
Combining with energy,
\[
  -gL\cos\theta = v_{0}^{2} - 2gL(1-\cos\theta)
  \;\Longrightarrow\;
  v_{0}^{2} = gL(2 - 3\cos\theta).
\]
The slack angle $\theta'$ therefore satisfies $\cos\theta' = (2 - v_{0}^{2}/gL)/3$.  The condition $\theta' > \pi/2$ means $\cos\theta' < 0$, i.e.\ $v_{0}^{2} > 2gL$.  Hence
\[
  \boxed{v_{0,\min} = \sqrt{2gL}}.
\]

\subsection*{Q2. Over the top}

The condition $\theta' > \pi$ means $\cos\theta' < -1$ requires $v_{0}^{2} > 5gL$, so the boundary is $\theta' = \pi$, i.e.\ $\cos\theta' = -1$.  Then
\[
  v_{0}^{2} = gL(2 + 3) = 5gL
  \;\Longrightarrow\;
  \boxed{v_{0,\min} = \sqrt{5gL}}.
\]
At this speed the string goes slack exactly at the top, $\theta' = \pi$.

\subsection*{Q3. The point of no return}

Substituting the energy relation $v^{2} = v_{0}^{2} - 2gL(1-\cos\theta)$ into the tension gives
\[
  T(\theta) = M\frac{v^{2}}{L} + Mg\cos\theta
            = \frac{M}{L}\,v_{0}^{2} + Mg\bigl(3\cos\theta - 2\bigr).
\]
Since $\cos\theta$ is minimised at $\theta = \pi$, the tension is minimised at the top:
\[
  T_{\min} = T(\pi) = \frac{M}{L}\bigl(v_{0}^{2} - 5gL\bigr).
\]

Therefore, if the string is still taut at $\theta = \pi$ (so $T(\pi) > 0$), then $T(\theta) \ge T(\pi) > 0$ for every $\theta$, and in particular for every $\theta > \pi$.  So $\theta^{*} = \pi$ has the required property.

It is also the \emph{minimum}: for any $\theta^{*} < \pi$, choose a trajectory whose slack angle satisfies $\theta^{*} < \theta' < \pi$ (such trajectories exist, from Q1 with $\cos\theta' = (2 - v_{0}^{2}/gL)/3$).  Then $T(\theta^{*}) > 0$ yet $T(\theta') = 0$, so the property fails.  Hence no $\theta^{*} < \pi$ works, and
\[
  \boxed{\theta^{*} = \pi}.
\]
Equivalently: the string remains taut for the whole revolution precisely when $v_{0}^{2} \ge 5gL$, and the top is the last point at which it can first go slack.

\end{solution}

\input qpd-end
